\documentclass[12pt]{article}
\usepackage[utf8]{inputenc}
\usepackage[T1]{fontenc}
\usepackage{hyperref}
\title{Dalecarlian Language Resources for K-12 Curriculum Development}
\author{Compiled by Claude Sonnet 5.5 (claude-sonnet-5-5) for Andr\'as Kornai}
\date{}
\begin{document}
\maketitle

\section{Introduction}
Dalecarlian (ISO 639-3 qer; Glottocode dale1238) is a North Germanic group of varieties in Dalarna, Sweden. The best-known variety is Elfdalian (\"Ovdalian, \"ovdalsk), in \"Alvdalen. The searches were run mainly on Elfdalian; other Dalecarlian varieties were not separately searched. No dedicated dictionary or complete school-material set turned up in the results.

\section{Primary Educational Resources}
The only teaching-oriented item found is \emph{A Grammar of Elfdalian} by Olof Lundgren and Yair Sapir (UCL Press), described as intended partly for those teaching the language in preschool, school or adult classes:
\begin{itemize}
\item \url{https://library.oapen.org/handle/20.500.12657/94747}
\item \url{https://press.uchicago.edu/ucp/books/book/distributed/G/bo238314327.html}
\end{itemize}
Reports indicate that Elfdalian is taught occasionally in \"Alvdalen schools and preschools, but no curriculum documents were found.

\section{Linguistic and Reference Materials}
\begin{itemize}
\item Wikipedia, \"Ovdalian: \url{https://en.wikipedia.org/wiki/%C3%96vdalian}
\item Omniglot, Elfdalian alphabet and pronunciation (standard orthography of 2005 by the Elfdalian Language Council, R\aa\dh djarum): \url{https://www.omniglot.com/writing/elfdalian.htm}
\end{itemize}

\section{Religious and Community Resources}
No religious texts were found in these searches. The language council R\aa\dh djarum and the association Ulum Dalska were mentioned as places to look next, but no URLs for them were seen.

\section{Digital and Audio Resources}
None found.

\section{Curriculum Development Context}
Policy overview: \url{https://wiki.mercator-research.eu/_export/raw/languages:elfdalian_in_sweden}. A news report describes a stipend for \"Alvdalen pupils who finish ninth grade and can show they speak the language: \url{https://www.thenewdaily.com.au/news/2015/05/18/sweden-ancient-elvish-language-dying}. A standard orthography and a modern grammar exist, giving a base for materials, but primers and textbooks would largely need to be created.
\end{document}